\documentclass[9pt]{article}

% --- Packing + layout ---
\usepackage{fancyhdr}
\usepackage[letterpaper,margin=0.3in,includehead,includefoot,headheight=14pt,headsep=6pt]{geometry}
\setlength{\columnsep}{0.15in}
\usepackage{microtype}
\usepackage{multicol}
\setlength{\parindent}{0pt}
\setlength{\parskip}{0pt}

% Tighten section spacing
\usepackage[compact,small]{titlesec}
\titlespacing*{\section}{0pt}{2pt}{1pt}
\titlespacing*{\subsection}{0pt}{2pt}{1pt}

% Tighten lists
\usepackage{enumitem}
\setlist{nosep,leftmargin=*}
\setlist[itemize]{noitemsep,topsep=0pt,parsep=0pt,partopsep=0pt,leftmargin=*}
\setlist[enumerate]{noitemsep,topsep=0pt,parsep=0pt,partopsep=0pt,leftmargin=*}

% Math spacing
\setlength{\abovedisplayskip}{2pt}
\setlength{\belowdisplayskip}{2pt}
\setlength{\abovedisplayshortskip}{1pt}
\setlength{\belowdisplayshortskip}{1pt}

% Math + links
\usepackage{amsmath,amssymb}
\usepackage[hidelinks]{hyperref}

% For including images
\usepackage{graphicx}
\usepackage{tikz}
\usetikzlibrary{calc}

% Callout boxes (kept for consistency even if unused)
\usepackage[most]{tcolorbox}
\tcbset{colback=white,colframe=black,boxrule=0.4pt,arc=2pt,left=6pt,right=6pt,top=4pt,bottom=4pt}
\newtcolorbox{notebox}{title=Note}
\newtcolorbox{solutionbox}{title=Solution}

% --- Fancy header/footer ---
\setlength{\headheight}{14pt}
\pagestyle{fancy}
\fancyhf{}
\fancyhead[L]{\textbf{Core Assessment 4 2026-02-10}}
\fancyfoot[C]{\thepage}
\renewcommand{\headrulewidth}{0.4pt}
\fancypagestyle{plain}{%
  \fancyhf{}
  \fancyhead[L]{\textbf{Core Assessment 5}}
  \fancyfoot[C]{\thepage}
  \renewcommand{\headrulewidth}{0.4pt}
}
\newcommand{\MinusOne}{\rulecite{Sign-flip: $a\equiv -b$}}
\newcommand{\Period}{\rulecite{Periodicity / Euler–Carmichael}}
\newcommand{\CRT}{\rulecite{CRT}}
\newcommand{\Geo}{\rulecite{Geometric grouping}}
\newcommand{\rulecite}[1]{\textit{[#1]}}
\newcommand{\CDL}{\rulecite{CDL}}
\newcommand{\FD}{\rulecite{FD}}
\newcommand{\Residues}{\rulecite{Residues}}
\newcommand{\Parity}{\rulecite{Parity}}
\newcommand{\Factor}{\rulecite{Factor}}
\newcommand{\EL}{\rulecite{Euclid's Lemma}}
\newcommand{\LinComb}{\rulecite{Linear Combinations}}
\newtheorem{problem}{Problem}
\newenvironment{solution}{\par\noindent\textbf{Solution.}\ }{\hfill$\square$\par}

\begin{document}
\small
\begin{problem} \qquad
\newline
\textbf{Definition (Factorial).}
For any integer $n \ge 0$, the factorial of $n$ is
\[
n! = n\,(n-1)\,(n-2)\cdots 2\cdot 1,
\qquad \text{with the convention } 0! := 1.
\]

\medskip
\noindent
\textbf{Definition (Binomial Coefficient).}
For integers $n \ge 0$ and $0 \le k \le n$,
\[
\binom{n}{k} \;=\; \frac{n!}{k!\,(n-k)!}.
\]

\bigskip
\textbf{Tasks.}
\begin{enumerate}[label=\textbf{(\alph*)}]
  \item Prove \textbf{Pascal's identity}:
  \[
  \binom{n}{k}=\binom{n-1}{k-1}+\binom{n-1}{k}
  \qquad (n\ge 1,\ 1\le k\le n-1).
  \]
  You may use either a combinatorial argument or algebraic manipulation from the factorial definition.

  \item Prove by \textbf{mathematical induction} on $n$ that
  \[
  \sum_{k=0}^{n}\binom{n}{k}=2^n
  \qquad \text{for all } n\ge 0.
  \]
  (\emph{Hint:} Use Pascal's identity in the inductive step.)
\end{enumerate}
\end{problem}

\begin{solution}\qquad

\textbf{(a) Pascal's identity.}
\vfill

\textbf{(b) Induction for $\displaystyle \sum_{k=0}^{n}\binom{n}{k}=2^n$.}

\smallskip
\emph{Base case ($n=0$).}

\vfill
\emph{Inductive step.}
\end{solution}

\vfill
\vfill
\end{document}